\subsection{RQ5: Practitioner Perceptions of Goal Violations}
\label{ss:sec:rq5}

\reva{
    11 of 13 respondents (84.6\%) found the goal violation examples useful to their workflow, scoring their utility with a median of five on a five-point scale.
    Five respondents highlighted its utility for guiding further review, with one response putting it as ``\textit{a triage and discovery layer, but not \ldots a final judgment.}''
    Three respondents highlighted the concrete goal-violating transaction sequences as useful for understanding the issue.
}

\reva{
    All respondents rated the end-user behaviors %with end-user-facing behaviors 
    that goal violations uncover as relevant to security, with a median of five on a five-point scale.  This is despite the fact that the provided examples are deliberate features of the USDT contract. 
    One respondent highlighted the context-dependent nature of whether a goal violation was concerning, noting that ``\textit{these violations on this particular contract don't mean much, but I can see these kinds of flags being useful when buying novel assets}''.
    All respondents believed that most end-users were unaware of administrative risks such as pausing transfers or blacklisting addresses, with one respondent commenting that ``\textit{most users assume ‘decentralized’ and miss centralized admin controls entirely}''. Overall,  this result demonstrates the usefulness of discovering end-user goal violations for security practitioners and  \spoilsport{}'s utility in aiding  smart contract auditing and validation.  %'s,  findings (end-user goal violations) . 
}

\reva{\begin{results}
    11 of 13 respondents (84.6\%) found the goal violation examples useful to their workflow.
\end{results}
}
